% !TEX root = TFM-ML-01-memoria-referencia.tex
\section{Evaluación}\label{sec:evaluacion}\label{cap:replicabilidad}

%% ══════════════════════════════════════════════════════
%% §8.1 — PLIEGO DE CONDICIONES DEL SERVICIO
%% Estado: REDACTADO
%% Trazabilidad: insights (cap05) → QFD (cap06) → aquí → matriz §8.2
%% ══════════════════════════════════════════════════════

\subsection{Pliego de condiciones del servicio de movilidad nocturna}
\label{sec:pliego}

El presente pliego de condiciones establece el marco de requerimientos
que debe satisfacer el concepto de servicio de movilidad nocturna para
trabajadores sanitarios del corredor \acs{CHUF}--Ferrol. Actúa como
documento de síntesis estratégica: formaliza las exigencias derivadas
de la investigación de usuario y los criterios de diseño, y establece
los ejes con los que se evaluará la propuesta en la sección~\ref{sec:matriz-evidencia}.

%% ── 8.1.1 OBJETO DEL ENCARGO ──────────────────────────

\subsubsection{Objeto del encargo}
\label{sec:pliego-objeto}

El presente pliego tiene por objeto establecer los requerimientos para
el diseño de un \textbf{sistema de movilidad nocturna} orientado a
resolver el vacío de cobertura de transporte público para trabajadores
sanitarios a turnos en el corredor \acs{CHUF}--Ferrol, en la franja
horaria de 22:00 a 08:00\,h.

La arquitectura del servicio (puntos de embarque, modo de captación
de demanda y tecnología habilitadora) se determina en el
capítulo~\ref{cap:adn-concepto} tras la evaluación de candidatos
(§\ref{sec:evaluacion-candidatos}). Los requerimientos que siguen son
independientes de esa selección.

El encargo se formula en el marco del Trabajo Fin de Máster en
Ingeniería en Diseño Industrial (EUDI, Universidade da Coruña),
Especialidad Automoción, Movilidad y Transporte. El alumno adopta
el rol de Promotor--Diseñador, estableciendo las reglas desde la
investigación empírica y evaluando la propuesta resultante contra
dichas reglas.

%% ── 8.1.2 DOCUMENTACIÓN DE REFERENCIA ────────────────

\subsubsection{Documentación de referencia obligatoria}
\label{sec:pliego-documentacion}

La propuesta debe fundamentarse estrictamente en los siguientes
análisis previos, desarrollados en los capítulos anteriores de
esta memoria:

\begin{itemize}
  \item Análisis \acs{PESTEL} del macroentorno europeo de movilidad
        nocturna (§\,\ref{sec:pestel}).
  \item Caso de estudio Ferrol y mapa de stakeholders
        (§\,\ref{sec:caso-ferrol}--\ref{sec:stakeholders}).
  \item Benchmarking de pilotos europeos de movilidad autónoma en
        entorno hospitalario (§\,\ref{sec:benchmarking}).
  \item Perfiles de usuario y vulnerabilidades detectadas, construidos
        sobre encuesta primaria $n = 167$
        (§\,\ref{sec:perfil-usuario}).
  \item Los cinco insights de diseño validados, con especial atención
        a la fatiga al volante post-turno y la brecha de género en el
        riesgo nocturno (§\,\ref{sec:insights-diseno}).
  \item Criterios de diseño sintetizados mediante \acs{QFD} y
        priorización \acs{MoSCoW} (§\,\ref{sec:qfd}--\ref{sec:moscow}).
\end{itemize}

%% ── 8.1.3 PERFIL DEL USUARIO OBJETIVO ────────────────

\subsubsection{Perfil del usuario objetivo}
\label{sec:pliego-usuario}

El \textbf{perfil primario de diseño} es el personal de enfermería y
\acs{TCAE} con turno nocturno fijo (22:00--08:00\,h), mujer en el
91{,}0\,\% de los casos, con patrón de salida sincronizado por planta
y mayor concentración numérica en la muestra
\parencite{LouridoRegueira2026}.

El \textbf{perfil secundario} corresponde al personal médico en régimen
de guardia de atención continuada (15:00--08:00\,h o 24\,h en festivo),
con horarios variables y menor volumen en la muestra.

El servicio se diseña para el perfil primario. La flota y los nodos
\acs{PR} se dimensionan a partir de este perfil. El servicio
\textbf{no} opera en modalidad puerta a puerta: conecta nodos de
aparcamiento periférico con el hospital, no domicilios con el hospital.

%% ── 8.1.4 ESPECIFICACIONES DEL SERVICIO ──────────────

\subsubsection{Especificaciones del servicio}
\label{sec:pliego-especificaciones}

Las especificaciones siguientes se derivan directamente de los insights
de diseño (cap.~\ref{cap:insights}). Ninguna especificación puede
formularse sin que un insight o dato verificado la respalde.

\begin{description}
  \item[E1 — Operación en franja 22:00--08:00\,h.]
        El servicio opera exclusivamente durante la franja de mayor
        riesgo de fatiga al volante post-turno, sin cobertura diurna.
        \emph{Fuente: Insight~1 + Insight~3.}

  \item[E2 — Nodos de embarque en zonas de aparcamiento libre o seguro.]
        Los puntos de acceso al servicio se ubican en zonas donde el
        trabajador puede dejar el vehículo sin coste ni demora,
        eliminando la necesidad de aparcar en el recinto hospitalario.
        \emph{Fuente: Insight~2.}

  \item[E3 — Sincronización con cambios de turno.]
        Los horarios de servicio se vinculan a los cambios de turno
        principales (22:00, 00:00 y 07:30\,h), no a frecuencias
        genéricas de red urbana.
        \emph{Fuente: Insight~3 + perfil primario.}

  \item[E4 — Modelo punto a punto: nodo P\&R → hospital.]
        El servicio conecta nodos periféricos con el \acs{CHUF}, no
        actúa como extensión de línea urbana ni como servicio
        puerta a puerta.
        \emph{Fuente: Insight~4.}

  \item[E5 — Seguridad percibida para usuarias que viajan solas en
        franja nocturna.]
        El diseño del vehículo, los nodos de espera y los touchpoints
        garantizan la accesibilidad y la seguridad percibida para la
        usuaria principal, que viaja sola de noche en el 91{,}0\,\%
        de los casos.
        \emph{Métrica: puntuación de seguridad percibida
        $\geq$\,4/5 en encuesta post-servicio a usuarias que viajan
        solas; iluminación mínima de 50~lux en todos los nodos
        de espera.}
        \emph{Fuente: Insight~5.}

  \item[E6 — Mínimo esfuerzo cognitivo post-turno.]
        La reserva, la espera y el embarque deben requerir el mínimo
        esfuerzo cognitivo posible de una persona en estado de fatiga
        tras 10--12\,h de turno nocturno.
        \emph{Métrica: acceso al servicio en un máximo de
        3~interacciones desde la aplicación o el panel del nodo;
        tiempo de activación $\leq$\,2~minutos desde la salida
        del turno.}
        \emph{Fuente: Insight~1.}

  \item[E7 — Replicabilidad a corredores análogos.]
        El modelo de servicio debe ser transferible a ciudades europeas
        de 50.000--150.000 habitantes con hospital periférico y vacío
        de transporte nocturno.
        \emph{Fuente: argumento de replicabilidad (cap.~\ref{cap:replicabilidad}).}
\end{description}

%% ── 8.1.5 CRITERIOS DE EVALUACIÓN PONDERADOS ─────────

\subsubsection{Criterios de evaluación de la propuesta}
\label{sec:pliego-criterios}

La propuesta de servicio resultante se evalúa contra tres ejes
estratégicos. Cada eje tiene asignado un peso porcentual que
refleja la prioridad metodológica derivada de la investigación
de usuario. La evaluación detallada se desarrolla en la matriz
de afirmaciones--evidencia (§\,\ref{sec:matriz-evidencia}).

\paragraph{CE1 — Seguridad percibida y experiencia de usuario (40\,\%)}

Evalúa en qué medida el diseño responde a las fricciones emocionales
documentadas: reducción del riesgo de fatiga al volante post-turno y
garantía de seguridad percibida para usuarias que viajan solas en
franja nocturna. Abarca el diseño del espacio interior, los nodos de
espera y la interacción con el sistema.
\emph{Especificaciones evaluadas: E5, E6.}

\paragraph{CE2 — Eficiencia logística y diferencial frente al modelo actual (30\,\%)}

Evalúa la coherencia operativa del servicio: sincronización con los
cambios de turno, eliminación del coste oculto del aparcamiento y
diferencial respecto al transporte público existente en la franja
nocturna.
\emph{Especificaciones evaluadas: E1, E2, E3, E4.}

\paragraph{CE3 — Viabilidad técnica, operativa y replicabilidad (30\,\%)}

Evalúa la factibilidad de operación del servicio en el corredor
\acs{CHUF}--Ferrol, su integración en el ecosistema de actores
(\acs{SERGAS}, Concello, operadores de transporte) y la
transferibilidad del modelo a corredores análogos en Europa.
\emph{Especificaciones evaluadas: E7 + compatibilidad stakeholders.}

%% ── 8.1.6 ENTREGABLES DEL PROYECTO ───────────────────

\subsubsection{Entregables del proyecto}
\label{sec:pliego-entregables}

Se requieren los siguientes entregables de diseño de servicio,
desarrollados en el cap.~\ref{cap:concepto}:

\begin{itemize}
  \item \emph{Service Blueprint} completo de la operation del servicio.
  \item \emph{Journey Map} de la usuaria primaria (enfermera/\acs{TCAE},
        turno 22:00--08:00\,h).
  \item Diseño de \emph{touchpoints} del servicio (reserva, espera,
        embarque y desembarque).
  \item Principios de diseño del espacio interior y de la interfaz
        hombre-máquina (\acs{HMI}).
  \item Modelo de negocio conceptual (\emph{Business Model Canvas}).
\end{itemize}

%% ── 8.1.7 MARCO NORMATIVO Y CONDICIONES DEL SERVICIO ─

\subsubsection{Marco normativo y condiciones del servicio}
\label{sec:pliego-normativa}

El diseño del servicio debe ser compatible con el marco normativo
vigente en materia de transporte de viajeros, accesibilidad y
protección de datos:

\begin{itemize}
  \item \textbf{Reglamento (UE) n.\textsuperscript{o} 181/2011} del
        Parlamento Europeo y del Consejo, de 16 de febrero de 2011,
        sobre derechos de los viajeros de autobús y autocar. Aplica
        en materia de no discriminación, accesibilidad para personas
        con discapacidad o movilidad reducida, e información al
        usuario.
  \item \textbf{Real Decreto 1544/2007}, de 23 de noviembre, por el
        que se regulan las condiciones básicas de accesibilidad y no
        discriminación para el acceso y utilización de los modos de
        transporte para personas con discapacidad (modificado por
        RD\,537/2019 y RD\,193/2023).
  \item \textbf{Reglamento Delegado (UE) 2024/490} de la Comisión,
        de 29 de noviembre de 2023, sobre suministro de servicios de
        información sobre desplazamientos multimodales. Cubre
        explícitamente los servicios de «autobús lanzadera» y
        «transporte público a la demanda».
  \item \textbf{Directiva (UE) 2023/2661} del Parlamento Europeo
        y del Consejo, de 22 de noviembre de 2023, que modifica
        la Directiva 2010/40/UE sobre Sistemas de Transporte
        Inteligentes (STI), marco legal aplicable a la movilidad
        automatizada.
  \item \textbf{Recomendación (UE) 2023/550} de la Comisión, de
        8 de marzo de 2023, sobre planes de movilidad urbana
        sostenible. Establece que los municipios incluyan movilidad
        compartida nocturna en sus planes de movilidad.
  \item \textbf{Reglamento (UE) 2016/679} (RGPD), sobre protección
        de datos de carácter personal. Aplica al tratamiento de datos
        de reserva y uso del servicio por parte de los usuarios.
\end{itemize}

%% ── 8.1.8 PLAZOS ──────────────────────────────────────

\subsubsection{Plazos}
\label{sec:pliego-plazos}

El presente pliego regula el horizonte de diseño del proyecto, no el
horizonte de implementación del servicio. Este último se declara en
el alcance del \acs{TFM} (§\,\ref{sec:alcance}).

El desarrollo completo del concepto de servicio —desde el cierre de
los criterios de diseño hasta la entrega de los entregables listados
en §\,\ref{sec:pliego-entregables}— se enmarca en el calendario del
\acs{TFM}, con fecha de entrega en junio de 2026. La estimación de
horas de desarrollo se recoge en el capítulo de presupuesto
(cap.~\ref{cap:presupuesto}).

\subsection{DAFO del servicio propuesto}\label{sec:dafo}

El análisis DAFO evalúa el concepto de servicio propuesto desde
dos dimensiones: la interna (fortalezas y debilidades del propio
sistema de movilidad diseñado) y la externa (oportunidades y
amenazas del contexto en el que opera el corredor CHUF--Ferrol).

\subsubsection{Fortalezas}
\label{sec:dafo-fortalezas}

\begin{itemize}

	\item \textbf{Demanda empíricamente documentada.} El servicio
	parte de una encuesta primaria (n\,=\,167) que cuantifica la
	necesidad con precisión: frecuencia, perfil de usuario, barreras
	actuales y demanda espontánea de solución
	\parencite{LouridoRegueira2026}. Ningún competidor actual en
	el corredor tiene este nivel de conocimiento del usuario.

	\item \textbf{Posición diferencial sin competencia directa.}
	Ningún operador existente (Tranvías de Ferrol, Monbus, Renfe
	Cercanías C-1 ni Cabify) ofrece servicio específico para
	personal sanitario en la franja 22:00--08:00
	\parencite{Alsaferrol2024horarios,RenfeCercaniasFerrol2026}.

	\item \textbf{Precedente europeo operativo.} El modelo P+R
	más lanzadera autónoma está en funcionamiento real en el
	campus hospitalario Belle-Idée de Ginebra y en el CHU de
	Rennes \parencite{AVENUE2023,SHOW2024Legacy}. El servicio no es
	una hipótesis: es una arquitectura probada en contexto
	hospitalario europeo.

	\item \textbf{Argumento de replicabilidad.} Ferrol encaja
	en el perfil de ciudad media europea (50.000--150.000 hab.)
	con hospital periférico y vacío de transporte nocturno
	\parencite{URBACTNight2024}. El modelo diseñado es
	transferible a corredores análogos, lo que amplía su
	valor estratégico más allá del caso local.

\end{itemize}

\subsubsection{Debilidades}
\label{sec:dafo-debilidades}

\begin{itemize}

	\item \textbf{Brecha tecnológica nocturna.} Ningún piloto
	documentado de vehículo autónomo SAE~L4 opera en franja
	nocturna plena a 2026 \parencite{SHOW2024Legacy}. La iluminación
	deficiente interrumpe la operación de los sensores. El
	horizonte 2035 de este TFM cubre esta brecha, pero introduce
	incertidumbre sobre la madurez tecnológica real en la
	fecha de implantación.

	\item \textbf{Infraestructura P+R inexistente.} Los nodos
	de aparcamiento disuasorio en los accesos al CHUF no
	existen actualmente. Su construcción requiere inversión
	pública previa e implica coordinación entre el Concello
	de Ferrol, la Xunta y el SERGAS.

	\item \textbf{Dependencia de un único corredor.} El modelo
	está diseñado para el corredor CHUF--Ferrol. Su eficacia
	depende de que la demanda esté suficientemente concentrada
	en ese eje. Si la dispersión residencial de los trabajadores
	es mayor de lo estimado, la rentabilidad operativa se reduce.

	\item \textbf{Ausencia de modelo de negocio cerrado.}
	El TFM propone una exploración de valor conceptual, no un
	plan financiero. La sostenibilidad económica del servicio
	requiere validación con actores institucionales que
	quedan fuera del alcance de este trabajo.

\end{itemize}

\subsubsection{Oportunidades}
\label{sec:dafo-oportunidades}

\begin{itemize}

	\item \textbf{Marco regulatorio europeo favorable.}
	El Reglamento (UE) 2026/481 habilita el despliegue a
	gran escala de vehículos totalmente automatizados
	\parencite{RegUE20262026481}. La ventana regulatoria
	para 2035 es la más abierta de la historia reciente
	en movilidad autónoma.

	\item \textbf{SERGAS como financiador potencial.}
	El sistema sanitario gallego tiene interés directo en
	la retención de personal nocturno \parencite{SERGAS2024memoria}.
	Un servicio de transporte específico puede encuadrarse
	como beneficio social en el convenio colectivo o como
	gasto de infraestructura sanitaria.

	\item \textbf{Agenda europea de movilidad sostenible.}
	La Estrategia de Movilidad Sostenible e Inteligente
	de la Comisión Europea \parencite{CE_Movilidad2020}
	y los fondos de cohesión territorial crean un contexto
	favorable para proyectos de transporte nocturno en
	ciudades medias.

	\item \textbf{Replicabilidad como factor de escala.}
	Si el modelo funciona en Ferrol, puede transferirse a
	ciudades análogas europeas \parencite{URBACTNight2024},
	lo que abre la posibilidad de financiación supramunicipal
	o de consorcio de ciudades.

\end{itemize}

\subsubsection{Amenazas}
\label{sec:dafo-amenazas}

\begin{itemize}

	\item \textbf{Dependencia de voluntad institucional.}
	El servicio requiere la coordinación activa de al menos
	tres administraciones (SERGAS, Xunta, Concello de Ferrol).
	La fragmentación competencial es el principal riesgo
	de bloqueo no tecnológico.

	\item \textbf{Hábito consolidado del vehículo privado.}
	El modo de transporte dominante en la muestra es el
	vehículo propio \parencite{LouridoRegueira2026}. El
	cambio de comportamiento modal requiere incentivos
	sostenidos y tiempo de adopción que el servicio debe
	anticipar en su diseño.

	\item \textbf{Riesgo de lock-in tecnológico.} Los
	pilotos AVENUE documentan dependencia total de las
	APIs propietarias de los fabricantes de vehículos
	autónomos \parencite{AVENUE2023}. La elección del
	proveedor tecnológico tiene implicaciones de
	sostenibilidad a largo plazo que superan el alcance
	de este TFM.

	\item \textbf{Volatilidad del marco tecnológico.}
	El horizonte 2035 introduce incertidumbre sobre qué
	tecnología estará disponible y a qué coste. Una
	apuesta excesivamente específica en la solución
	técnica puede quedar obsoleta antes de la implantación.

\end{itemize}

\begin{figure}[htbp]
	\centering
	\fbox{\rule{0pt}{7cm}\rule{0.95\textwidth}{0pt}}
	% Sustituir por: \includegraphics[width=\textwidth]{../03-graficas/TFM-ML-GFX-dafo.pdf}
	\caption{Análisis DAFO del servicio de movilidad nocturna
		propuesto para el corredor CHUF--Ferrol. Elaboración propia.}
	\label{fig:dafo}
\end{figure}




\subsection{Matriz de afirmaciones-evidencia}
\label{sec:matriz-evidencia}

La matriz evalúa las afirmaciones centrales del proyecto,
clasificando cada una según su estado de validación
(Supported / Contested / Undetermined) y el nivel de confianza
de la evidencia que la respalda (Alta / Media / Baja). Una
afirmación es Supported solo si existe fuente primaria
verificable o precedente operativo documentado que la sustente
de forma directa~\parencite{LouridoRegueira2026}.

\begin{table}[htbp]
\centering
\caption{Matriz de afirmaciones-evidencia del proyecto.
  Elaboración propia.}
\label{tab:matriz-evidencia}
\small
\begin{tabularx}{\textwidth}{p{5.2cm} c c p{4.8cm}}
\toprule
\textbf{Afirmación} &
\textbf{Estado} &
\textbf{Nivel} &
\textbf{Evidencia} \\
\midrule

Existe vacío de transporte público en el corredor
\acs{CHUF}--Ferrol en franja 22:00--06:00 &
Supported & Alta &
Horarios Alsa Ferrol y Renfe Cercanías;
39/167 encuestados sin servicio
\parencite{Alsaferrol2024horarios,
RenfeCercaniasFerrol2026,
LouridoRegueira2026} \\
\addlinespace

El 52,7\,\% de la muestra conduce con sueño o fatiga
varias veces a la semana &
Supported & Alta &
Encuesta primaria n=167; índice de gravedad
60,2/100~\parencite{LouridoRegueira2026} \\
\addlinespace

La brecha de género en riesgo de fatiga al volante
justifica un enfoque de diseño diferenciado &
Supported & Alta &
90,8\,\% mujeres vs. 53,8\,\% hombres en muestra
91\,\% femenina~\parencite{LouridoRegueira2026} \\
\addlinespace

El modelo lanzadera en corredor hospitalario es viable
en contexto europeo &
Supported & Alta &
\acs{CHU} de~Rennes y Belle-Idée~(Ginebra) en operación
regular~\parencite{SHOW2024Legacy,AVENUE2023} \\
\addlinespace

La aceptación del shuttle autónomo entre usuarios de
transporte público es favorable &
Supported & Media &
Valoración media $>$7/9 Likert en pilotos SHOW; contexto
diurno, no nocturno~\parencite{SHOW2024Legacy} \\
\addlinespace

La operación \acs{SAE}~L4 nocturna es técnicamente viable
a 2026 &
Contested & Baja &
Pilotos SHOW documentan limitaciones de sensórica en baja
iluminación; brecha TRL~4--5 nocturno~\parencite{SHOW2024Legacy} \\
\addlinespace

La operación \acs{SAE}~L4 nocturna será viable en el
horizonte 2035 &
Undetermined & Media &
No existe piloto nocturno documentado a esta fecha; el
horizonte 2035 es una proyección razonada, no un dato
verificado~\parencite{KonstantasAVENUE2023lessons} \\
\addlinespace

El acceso por bono mensual es el modelo preferido por
el usuario &
Supported & Baja &
Preferencia expresada por Participante~1 (n=1); no
generalizable desde una sola entrevista
(Participante~1, mayo de 2026) \\
\addlinespace

El servicio es replicable en ciudades europeas de
50.000--150.000 hab. con hospital periférico &
Supported & Media &
Pilotos SHOW en Trikala (81k~hab.), Herford (67k~hab.)
y Linköping (160k~hab.)~\parencite{SHOW2024Legacy,
URBACTNight2024} \\
\addlinespace

El modelo reduce la dependencia del vehículo privado
en el corredor \acs{CHUF}--Ferrol &
Undetermined & Baja &
Afirmación prospectiva; no existe dato operativo local.
La reducción depende de la tasa de adopción real del
servicio, fuera del alcance de este trabajo \\

\bottomrule
\end{tabularx}
\end{table}

Las afirmaciones en estado Undetermined o con nivel Baja
no invalidan la propuesta: delimitan con precisión el alcance
del trabajo. Este proyecto formula un modelo de servicio
justificado desde la evidencia disponible a 2026; su validación
operativa requiere un piloto en el corredor real.

\subsection{Valor diferencial del concepto}
\label{sec:valor-diferencial}

El valor diferencial evalúa el posicionamiento estratégico del
servicio frente a las alternativas existentes en el corredor,
no su viabilidad técnica (resuelta en la matriz de evidencia)
ni sus riesgos operativos (recogidos en el \acs{DAFO}).

\subsubsection{Adecuación a los perfiles de usuario}

El servicio está diseñado para un segmento específico y
cuantificado: enfermeras y~\acs{TCAE} de turno nocturno fijo
(22:00--08:00), que representan el 77,9\,\% de la muestra y el
perfil de mayor concentración numérica en el
\acs{CHUF}~\parencite{LouridoRegueira2026}. Ninguna de las
alternativas actuales (autobús urbano, taxi, vehículo propio)
está diseñada para este perfil: el autobús no opera en la
franja; el taxi no es asequible a escala mensual; el vehículo
propio impone la carga de conducción que el servicio elimina.
El encaje no es parcial: el servicio resuelve los tres
principales problemas declarados por la muestra de forma
simultánea (fatiga al volante, aparcamiento y ausencia de
\acs{TP})~\parencite{LouridoRegueira2026}.

\subsubsection{Originalidad formal y simbólica}

El servicio reencuadra el trayecto nocturno: no es tiempo
muerto de desplazamiento, es tiempo de recuperación entre
turnos. Este reencuadre no es retórico; está respaldado por
el diseño del espacio interior (P1--P5, §\ref{sec:espacio-interior})
y por el modelo de acceso sin fricción cognitiva. La asociación
institucional con el \acs{SERGAS} o el \acs{CHUF} como socio
operativo da al servicio una identidad reconocible y
diferenciada del ride-hailing genérico: es un beneficio laboral,
no un sustituto del taxi.

\subsubsection{Claridad estratégica}

El canal de implantación es doble: acuerdo institucional
\acs{SERGAS}/\acs{CHUF} para la gobernanza operativa y
suscripción mensual individual para el acceso del usuario.
Este canal evita la dependencia de una sola parte y distribuye
la financiación entre institución y usuario. El servicio no
compite con el transporte público urbano porque ese transporte
no existe en la franja nocturna del corredor; compite con el
vehículo privado y el taxi, alternativas frente a las que
la ecuación de valor del servicio es favorable tanto en coste
como en seguridad percibida~\parencite{PlyushetvaBoussauw2020}.

\subsubsection{Proyección comercial}

La sostenibilidad económica del servicio descansa sobre tres
palancas complementarias: la suscripción mensual del usuario
(coste comparable al aparcamiento hospitalario de pago,
estimado en~80~€/mes~\parencite{LouridoRegueira2026}); la
cofinanciación institucional del \acs{SERGAS} en el marco de
las Obligaciones de Servicio Público~(\acs{OSP})
\parencite{UITP2024}; y la reducción estructural de costes
operativos a largo plazo derivada de la eliminación del
conductor en la solución \acs{SAE}~L4. La replicabilidad del
modelo a ciudades europeas de 50.000--150.000 hab. con
hospital periférico amplía el mercado potencial más allá
del caso Ferrol~\parencite{URBACTNight2024}.

\subsubsection{Encaje con las tendencias del sector}

El servicio se alinea con cuatro tendencias documentadas en
el análisis \acs{PESTEL} (§\ref{sec:pestel}): la agenda de
movilidad sostenible de la Comisión Europea~(Green Deal,
\acs{SUMP})~\parencite{CE_Movilidad2020}; la maduración
regulatoria del vehículo autónomo en la~\acs{UE}
(Reglamento~2026/481)~\parencite{RegUE20262026481}; la
visibilización de la movilidad nocturna como objeto de
política pública municipal~\parencite{URBACTNight2024}; y
la presión estructural sobre los sistemas sanitarios europeos
para retener personal nocturno en un contexto de escasez
de profesionales~\parencite{Eurofound2020Sectors}.

\subsection{Viabilidad y riesgos}
\label{sec:viabilidad}
El sistema actual genera una ineficiencia documentada: trabajadoras
de la misma planta toman taxis individuales en ventanas de 10--15~min
por la salida escalonada (Participante~1, entrevista semiestructurada,
mayo de 2026). Coste actual: aproximadamente 7~euros por trayecto y
persona. Con 3--5~trabajadoras de la misma planta compartiendo vehículo
en la misma ventana horaria, el coste per cápita cae a 1--2~euros por
trayecto~--- una reducción del 70--85\,\%.
Esta estimación es de orden de magnitud y procede de una única fuente
cualitativa (Participante~1, entrevista semiestructurada, mayo de~2026);
requiere contraste con las tarifas oficiales publicadas por los
operadores de taxi del municipio antes de emplearse como dato operativo.

\begin{pendiente}
\end{pendiente}
